\documentclass[11pt,a4paper]{article}
\usepackage[T1]{fontenc}
\usepackage{lmodern,amsmath,amssymb,amsthm,mathtools,booktabs,microtype}
\usepackage[margin=27mm]{geometry}
\usepackage[colorlinks=true,linkcolor=blue,citecolor=blue,urlcolor=blue]{hyperref}
\setlength{\emergencystretch}{2em}
\newtheorem{theorem}{Theorem}[section]
\newtheorem{lemma}[theorem]{Lemma}
\newtheorem{corollary}[theorem]{Corollary}
\newtheorem{proposition}[theorem]{Proposition}
\newcommand{\Q}{\mathbb Q}\newcommand{\Z}{\mathbb Z}
\DeclareMathOperator{\Tr}{Tr}\DeclareMathOperator{\Nm}{N}
\title{A finite-class obstruction for the native scalar source}
\author{}\date{16 November 2026}
\begin{document}\maketitle
\begin{abstract}
For every finite set of nontrivial rational square classes, including
negative classes, we construct an arithmetic progression of integer
inputs for which the two native general-source complements cannot both
belong to the specified classes. The obstruction holds for arbitrary
guarded rational auxiliary tuples and slopes, for all nonzero rational
level multipliers, and uniformly in the odd scalar exponent. The proof
uses three quartic root fields, Chebotarev and the source's exact
unit-weight alternative. It does not exclude the rational-square class
or give a lower bound on defining-polynomial variables.
\end{abstract}

\section{The literal native source}
Use the current unsquared general source of~\cite{UnitNorm}. In particular,
\begin{equation}\label{eq:scalar}
\begin{gathered}
 \beta(X)=216X^4+450459,\quad\ell=76249,\\
 r=\beta(\beta-1)(\beta-2401)(32\beta+343)(2\beta-1),\\
 C_L=17(\ell^{(L+1)/2}+1),\qquad
 \rho=\frac{102^2\ell^Lr}{(1+C_Lr)^2}\quad(L\ge3\text{ odd}).
\end{gathered}
\end{equation}
The scales and contact functions are
\begin{equation}\label{eq:contacts}
\begin{gathered}
 a_0=1,\quad a_1=1+\rho T^2,\quad t_j=a_j\lambda,\\
 \Theta=\{0,-1/2,-1,1/2399,-32/407\},\\
 \Theta_j=(4j+2+\Theta)\cup(4j+4+\Theta),\quad\gamma_j=-j-1,\\
 R_j=\frac{\prod_{\theta\in\Theta_j}(t_j-\theta)}{(t_j-\gamma_j)^{10}},
 \quad w_j=\rho R_j,\quad
 V_j=\Tr(\mathcal W_j^{-1}),\quad b_j=4w_j/V_j^2.
\end{gathered}
\end{equation}
The inherited guards require all displayed denominators, $w_j$ and
$V_j$ to be nonzero. The argument uses only~\eqref{eq:scalar}--
\eqref{eq:contacts}; it imposes no further restriction on the
elliptic auxiliary tuples.

Let $S$ be the original finite exceptional set. For $m\in\Q^\times$ put
\[
                         C_m(b)=17m^2(17m^2-b).
\]

\begin{theorem}\label{thm:main}
Let $\mathcal D$ be any nonempty finite set of nontrivial rational
square classes. There is an arithmetic progression of integer inputs
$X$, independent of $L$, such that no guarded rational tuple of
\eqref{eq:scalar}--\eqref{eq:contacts} satisfies
\begin{equation}\label{eq:conditions}
              d_jC_{m_j}(b_j)\in\Q^2\qquad(j=0,1)
\end{equation}
with $d_0,d_1\in\mathcal D$ and $m_0,m_1\in\Q^\times$.
This holds for every fixed odd $L\ge3$ and for arbitrary rational
slopes and auxiliary indices satisfying the inherited guards.
\end{theorem}

The class representatives need not be positive or supported on $S$.
The theorem is an obstruction for the specified native scalar, not
for other scalar input maps or all possible Diophantine constructions.
It does not depend on the general source's external arithmetic premises.

\begin{corollary}\label{cor:split}
Allow all positive square classes supported on $S$. For some integer
inputs, every successful two-complement tuple must use the rational-square
class at at least one detector.
\end{corollary}
\begin{proof}
Apply Theorem~\ref{thm:main} to the finite set of all nontrivial
positive classes supported on $S$.
\end{proof}
This does not force both complements to be squares and does not
disprove realization when the rational-square class is included.

\section{The local norm obstruction}
\begin{lemma}\label{lem:norm}
For nonzero rational $w,V,m,d$ and $b=4w/V^2$,
\[
 dC_m(b)\in\Q^2
 \quad\Longleftrightarrow\quad
 \exists Y\in\Q:\quad(mV)^2-dY^2=4w/17.
\]
If $p\ne2,17$, $d$ is a $p$-adic unit and a nonsquare modulo $p$,
and $v_p(w)$ is odd, this condition is impossible for every $m$.
\end{lemma}
\begin{proof}
The identity
\[
 C_m(b)=\left(\frac{17m}{V}\right)^2
               \left((mV)^2-\frac{4w}{17}\right)
\]
gives the equivalence: if $z^2=dC_m(b)$, take
$Y=zV/(17md)$, and reverse this substitution for the converse.

For $A,B\in\Q_p$ not both zero, nonsquareness of $d$ gives
\[
                   v_p(A^2-dB^2)=2\min\{v_p(A),v_p(B)\}.
\]
Unequal term orders cannot cancel; equal orders could cancel only if
$d$ were a residue square. Thus the left side of the norm equation
has even valuation, whereas its right side has odd valuation.
\end{proof}

\begin{lemma}\label{lem:unit}
Let $p$ avoid the fixed finite set of contact denominators, zero
residues and cross-set collisions. If $v_p(\rho)$ is odd, then
at least one of $R_0,R_1$ is a unit at $p$.
\end{lemma}
\begin{proof}
If $T\ne0$, the valuation of $\rho T^2$ is odd and hence nonzero.
Thus $a_1\equiv1\pmod p$ or $v_p(a_1)<0$; the first alternative
also covers $T=0$. An argument of positive or negative valuation
gives a unit in a balanced contact function: all contact and pole
constants are units, and the numerator and denominator have equal degree.

If $v_p(\lambda)\ne0$, this gives a unit $R_0$. If $\lambda$ is a
unit and $R_0$ is not, its residue is a contact or pole of the first
set. Either $a_1\lambda$ has negative valuation or it has the same
residue as $\lambda$. The second contact-and-pole set is disjoint
from the first, so $R_1$ is a unit in either case. This is the unit
part of the scalar-odd argument in~\cite{General,UnitNorm}.
\end{proof}

Consequently, at a prime inert in both selected classes and of odd
$\rho$-valuation, at least one of the two conditions
\eqref{eq:conditions} is impossible. This observation does not use
a prime-prepared slope.

\section{Three quartic root fields}
For $\theta\in\{0,1,2401\}$ let $\alpha_\theta$ be a root of
\[
             216x^4+450459-\theta,
 \qquad F_\theta=\Q(\alpha_\theta),\qquad
 q_\theta=(450459-\theta)/216.
\]

\begin{lemma}\label{lem:fields}
Each $F_\theta$ has degree four and has a unique quadratic subfield.
The three quadratic subfields are respectively
\[
 \Q(\sqrt{-300306}),\qquad
 \Q(\sqrt{-675687}),\qquad
 \Q(\sqrt{-672087}).
\]
They are distinct imaginary fields.
\end{lemma}
\begin{proof}
The three quartics are Eisenstein at $50051,79,13$, respectively.
Their quadratic subfields $E_\theta=\Q(\alpha_\theta^2)$ follow from
\[
       q_0=50051/24,\quad q_1=225229/108,\quad q_{2401}=224029/108.
\]
These positive rationals are nonsquares.

Since $q_\theta$ is a nonsquare, $E_\theta\ne\Q(i)$.
The field $F_\theta$ does not contain $i$. Otherwise
$F_\theta=E_\theta(i)$, and writing a square root of
$\alpha_\theta^2$ as $a+bi$ with $a,b\in E_\theta$ shows that
$\alpha_\theta^2$ or its negative is a square in $E_\theta$.
Both alternatives contradict
$\Nm_{E_\theta/\Q}(\alpha_\theta^2)=q_\theta\notin\Q^2$.
Thus $F_\theta$ is not Galois, since the other roots include
$i\alpha_\theta$. Two distinct quadratic subfields would make a
quartic field biquadratic and Galois, proving uniqueness. The displayed
negative squarefree radicands are distinct.
\end{proof}

\begin{lemma}\label{lem:chebotarev}
For any two nontrivial rational square classes $d_0,d_1$, there are
a contact $\theta\in\{0,1,2401\}$ and infinitely many rational primes
$p$ such that $\beta(x)-\theta$ has a simple root modulo $p$ and
both $d_0,d_1$ are nonsquares modulo $p$. These primes may avoid
any prescribed finite set.
\end{lemma}
\begin{proof}
The field $M=\Q(\sqrt{d_0},\sqrt{d_1})$ has at most two imaginary
quadratic subfields: if two square classes are negative, their product
is positive. Choose $\theta$ so that $E_\theta$ is not contained in
$M$. Lemma~\ref{lem:fields} gives $F_\theta\cap M=\Q$: a quadratic
intersection would be $E_\theta$, and a quartic intersection would
make the non-Galois field $F_\theta$ equal to the Galois field $M$.

Let $N$ be the Galois closure of $F_\theta M$. The restriction of
$\operatorname{Gal}(N/F_\theta)$ to $M$ is all of
$\operatorname{Gal}(M/\Q)$, since the fixed field of its image is
$M\cap F_\theta$. There is an element of $\operatorname{Gal}(M/\Q)$
acting by minus one on both chosen square roots: the classes are
either equal or independent. Lift it to an element fixing
$\alpha_\theta$.

Chebotarev's theorem~\cite[Theorem~8.31]{Milne} gives infinitely many
unramified primes with this Frobenius conjugacy class. Each conjugate
fixes a root of the quartic, while its actions on the quadratic fields
remain unchanged. The Frobenius factorization correspondence
\cite[Corollary~8.22]{Milne} gives a root modulo $p$ and inertness
in both quadratic fields. Excluding discriminant and leading-coefficient
primes makes the root simple. Removing any further finite set preserves
infinitude.
\end{proof}

\section{Uniform input progressions}
\begin{proof}[Proof of Theorem~\ref{thm:main}]
For each unordered pair of elements of $\mathcal D$, allowing repetition,
choose a different prime and a contact as in Lemma~\ref{lem:chebotarev}.
Exclude $S$, $\ell$, all representative numerator and denominator
primes, the bad contact primes of Lemma~\ref{lem:unit}, and all primes
dividing the nonzero differences among the five finite roots of $r$.

At the selected simple root modulo $p$, choose a lift $x_p$ modulo
$p^2$ with
\[
                        v_p(\beta(x_p)-\theta)=1.
\]
All but one of the $p$ lifts have this property. The Chinese remainder
theorem gives an arithmetic progression of integer $X$ satisfying all
these congruences. At every selected prime exactly one factor of $r$
has valuation one and the other factors are units. Hence
\[
                            v_p(r)=v_p(\rho)=1.
\]
The denominator $1+C_Lr$ is congruent to one modulo $p$ for every
odd $L$; the input progression is therefore independent of $L$.

If a guarded tuple chose $d_0,d_1\in\mathcal D$, use the prime
assigned to that pair. Lemma~\ref{lem:unit} makes one $R_j$ a unit,
so $v_p(w_j)=1$. Both class representatives are nonsquares there,
and Lemma~\ref{lem:norm} contradicts the corresponding complement
condition, for every rational nonzero level $m_j$.
\end{proof}

\section{Exact examples and verification}
For $\mathcal D=\{409,433\}$ a single prime suffices:
\[
                           p=76511,\qquad X\equiv31907\pmod{p^2}.
\]
Both classes are nonsquares modulo $p$. For $g(x)=24x^4+50051$,
\[
 g(31907)/p\equiv4411\pmod p,\qquad
 g'(31907)\equiv10372\pmod p.
\]
The two contact-and-pole sets are nonzero and disjoint modulo $p$.
Thus the integer input $31907$ itself rules out both detectors using
these two classes, uniformly in $L$, even at arbitrary guarded slopes.

For classes $409$, $433$ and $177097$, use primes
$76511$, $76607$ and $76403$. Their inert-class pairs cover every
unordered pair, including repetition. It gives the progression
\[
 X\equiv r_0\pmod{M},\qquad
 \begin{aligned}
 r_0&=183039969870449349691519947271,\\
 M&=200542068630790781548849843561.
 \end{aligned}
\]
No single prime can be inert in all three multiplicatively related
classes; the pairwise construction is necessary for this certificate.

The standard-library program
\begin{center}\small
\nolinkurl{negative_results/computation/finite_nonsquare_classes_v47/check.py}
\end{center}
verifies the three quartic-field numerical certificates, the explicit
prime and contact computations, the exact valuation lifts, the CRT
progression and the source hashes. Chebotarev proves the infinite
statement for arbitrary finite class sets; finite computation is used
only for the stated explicit examples.

The class $1$ remains outside the obstruction. The theorem also does
not apply after replacing the scalar input polynomial or contact weight
by a different construction. In particular, it does not contradict a
separately proved norm-compatible scalar replacement.

\begin{thebibliography}{9}
\bibitem{UnitNorm}
\emph{A single norm determinant for the 25- and 13-variable constructions},
project manuscript, 14 November 2026,
\nolinkurl{papers/finite_algebra/unit_norm_20261114_v28.tex}.
\bibitem{General}
\emph{A defining polynomial in twenty-five variables: combined scalar
and comparison conditions}, project source manuscript,
\nolinkurl{papers/few_variables25/few_variables25_merged_20261110_v19.tex}.
\bibitem{Milne}
J.~S. Milne, \emph{Algebraic Number Theory}, version~3.08,
19 July 2020, Corollary~8.22 and Theorem~8.31,
\url{https://www.jmilne.org/math/CourseNotes/ANT.pdf}.
\end{thebibliography}
\end{document}
